\chapter{环论}

% ===================== 第一节 =====================
\section{环、交换环、除环与域}

\begin{definition}[交换环与单位元]
一个环 $R$ 叫做一个\textbf{交换环}，是指对于 $R$ 的任意两个元 $a$，$b$ 来说，都有
\[
ab=ba.
\]
在一个交换环里，对于任意正整数 $n$ 以及环的任意两个元 $a$，$b$ 来说，都有
\[
a^n b^n=(ab)^n.
\]

\textbf{单位元} \quad 在群论里我们已经看到了单位元的重要性。在环的定义里我们没有要求一个环要有一个对于乘法来说的单位元。但一个环假如有这样一个元，我们也可以想象，这个元也会占一个重要的地位。

\textbf{定义} \quad 一个环 $R$ 的一个元 $e$ 叫做一个\textbf{单位元}，是指对于 $R$ 的任意元 $a$ 来说，都有
\[
ea=ae=a.
\]
一般地，一个环未必有一个单位元。
\end{definition}

\begin{note}
交换环是给乘法补上了交换律，$n\times n$ 实矩阵是一个非交换环。
\end{note}

\begin{note}[环定义2 1]
环两种运算，加法有交换群的所有性质，乘法有幺半群的所有性质 \\
加法可以有交换率有结合律有逆元，乘法只有结合律不需要有逆元，有逆元叫单位除了0都有单位就是除环 \\
环的乘法对加法有分配律
\end{note}

\begin{proposition}[命题 2.1：负元、$0$ 元的运算]
令 $(R,+,\cdot)$ 是一个环，而 $a,b,c\in R$，则
\begin{align}
a0&=0a=0\\
a(-b)&=(-a)b=-(ab)\\
(-a)(-b)&=ab
\end{align}

$0$ 是单位元，不写符号就是乘法。意思是：任何元素乘以加法单位元 $0$ 都等于 $0$，乘法和 $0$ 的关系建立在分配律基础上可证。$-b$ 是 $b$ 的加法逆元，不是"减 $b$"，整个式子说明：负数乘法的规律——乘法中取负，可以移到任何一个因子上；在环中不需要交换律，这条等式依然成立（可用分配律证明）。
\end{proposition}

\begin{proposition}[命题 2.2：零环，加法单位元与乘法单位元相等]
令 $(R,+,\cdot)$ 是一个环，则 $R=\{0\}$ 当且仅当 $0=1$。

有一个重要的环是零环，它是最平凡的环，即 $(0,+,\cdot)$。它只有一个元素，既是加法单位元也是乘法单位元，定义为 $00=0$，$0+0=0$。很容易检验这是一个环。我们要说明的一个结论是，一个环是零环当且仅当加法单位元等于乘法单位元，即 $0=1$。
\end{proposition}

\begin{definition}[定义 2.2：交换环]
令 $(R,+,\cdot)$ 是一个环，我们说 $R$ 是一个交换环，当 $R$ 对乘法有交换律，即
\[
\forall a,b\in R,\;ab=ba.
\]
\end{definition}

\begin{note}[定义2 2交换环]
交换环是给乘法补上了交换率，$n\times n$实矩阵是一个非交换环
\end{note}

\begin{definition}[定义 2.3：单位与除环]
令 $(R,+,\cdot)$ 是一个环，则 $(R^\times,\cdot)$，是由 $R$ 中所有乘法可逆元构成的群。$R$ 中的乘法可逆元又被称为 $R$ 中的\textbf{单位}。

我们知道，对于一般的环（不是零环的环），$0\ne 1$。因此 $0$ 永远是没有乘法逆元的，这是因为对所有 $a\in R$，我们有 $0a=0\ne 1$。所以最好的情况可能就是 $R\setminus\{0\}=R^\times$，即所有非零元素都有乘法逆元。这样的环的性质是很好的，我们称它为一个\textbf{除环}。
\end{definition}

\begin{note}[定义2 4除环非零元素都是单位]
除环非零元素都有逆元。加法是阿贝尔群，乘法构成群而不是仅仅幺半群；乘法对加法有左右分配律。证明一个环是除环：只需证任意两个非零元相乘结果不为零。证明域：在除环基础上再证明乘法满足交换律。
\end{note}

\begin{definition}[环中逆元]
一个有单位元环的一个元 $b$ 叫做元 $a$ 的一个\textbf{逆元}，是指满足
\[
ba=ab=1.
\]
一个元 $a$ 最多只能有一个逆元。因为，假如 $a$ 有两个逆元 $b$ 和 $b'$，那么
\[
bab'=b(ab')=b1=b=(ba)b'=1b'=b'.
\]
\end{definition}

\begin{definition}[定义 2.5：域]
令 $(R,+,\cdot)$ 是一个环，我们称 $(R,+,\cdot)$ 是一个\textbf{域}，若它是一个交换的除环。
\end{definition}

\begin{note}[定义2 5域：交换的除环]
乘法对加法有分配律别忘了。
\end{note}

\begin{proposition}[命题 2.3：域的等价条件]
$(R,+,\cdot)$ 是一个域，当且仅当
\begin{align}
&(R,+)\text{ 是一个阿贝尔群}\\
&(R\setminus\{0\},\cdot)\text{ 是一个阿贝尔群}\\
&\text{乘法对加法有分配律}
\end{align}
\end{proposition}

\begin{proposition}[命题 2.5：环的直积还是环]
令 $((R_i,+_i,\cdot_i))_{i\in I}$ 是一族环，则它们的直积 $(\prod_{i\in I}R_i,+,\cdot)$ 还是一个环。

\textbf{证明：} 根据幺半群和群论对直积是保持的，我们立刻知道 $\prod_{i\in I}R_i$ 对加法构成群，对乘法构成幺半群。因此只须检验乘法对加法的左右分配律。根据对称性，我们只证明左分配律。

令 $(x_i)_{i\in I},(y_i)_{i\in I},(z_i)_{i\in I}\in\prod_{i\in I}R_i$，则
\begin{align}
(x_i)_{i\in I}\cdot((y_i)_{i\in I}+(z_i)_{i\in I})
&=(x_i\cdot_i(y_i+_i z_i))_{i\in I}\\
&=(x_i\cdot_i y_i+_i x_i\cdot_i z_i)_{i\in I}\\
&=(x_i)_{i\in I}\cdot(y_i)_{i\in I}+(x_i)_{i\in I}\cdot(z_i)_{i\in I}
\end{align}
因此，$(\prod_{i\in I}R_i,+,\cdot)$ 是一个环。验证分配律。
\end{proposition}

\begin{note}[定义2 8环的直积]
环的直积是一族环（多个环）的逐坐标运算，类比群的直积。
\end{note}

% ===================== 第二节 =====================
\section{子环与环同态}

\begin{definition}[子环与子除环]
一个环 $R$ 的一个子集 $S$ 叫做 $R$ 的一个\textbf{子环}，是指 $S$ 本身对于 $R$ 的代数运算来说作成一个环。

一个环 $R$ 的一个子集 $S$ 叫做 $R$ 的一个\textbf{子除环}，是指 $S$ 本身对于 $R$ 的代数运算来说作成一个除环。

同样，我们可以规定\textbf{子整环}、\textbf{子域}的概念。

一个环的非空子集 $S$ 作成一个子环的充要条件是
\[
a,b\in S\Rightarrow a-b\in S,\quad ab\in S.
\]
一个除环的一个子集 $S$ 作成一个子除环的充要条件是 (i) $S$ 包含一个不等于零的元；(ii) $a,b\in S\Rightarrow a-b\in S$，$a,b\in S,b\ne 0\Rightarrow ab^{-1}\in S$。
\end{definition}

\begin{note}[定义2 6子环]
子环继承父环的代数结构：加群是父群的子群，乘幺半群是父幺半群的子幺半群；有加法单位元和乘法单位元，对运算封闭。类比子群继承群运算。
\end{note}

\begin{lemma}[引理 2.1：子环等价条件]
令 $(R,+,\cdot)$ 是一个环，而 $S\subset R$，则 $S<R$ 当且仅当
\[
1\in S
\]
\[
\forall a,b\in S,\;a-b\in S,\quad ab\in S
\]
\textbf{证明：} 假如满足了这两个条件。首先 $0=1-1\in S$，而 $-a=0-a\in S$，$a+b=a-(-b)\in S$。这就证明了这是个子环（$a-b$ 蕴含着逆元）。

另一个方向是显然的。假如 $S$ 是子环，那么 $a-b=a+(-b)\in S$。有了子环，就可以利用用子集可以生成子幺半群。
\end{lemma}

\begin{lemma}[引理 2.2：理想是子环当且仅当理想是整个环]
$I$ 是 $R$ 的理想，$I$ 是 $R$ 的子环当且仅当 $I=R$：若 $1\in I$，则 $I$ 里元素左乘 $R$ 里任何元素都在 $I$ 里，故 $I=R$。
\end{lemma}

\begin{note}[引理2 2理想是子环当且仅当理想是整个环]
$I$ 是 $R$ 的理想；$I$ 是 $R$ 的子环当且仅当 $I=R$，即 $1\in I$ 时理想就是整个环。
\end{note}

\begin{lemma}[引理 2.3：理想可以从环同态中定义]
正如正规子群可以从群同态中定义（群同态的核），理想也可以从环同态中定义（环同态的核）。

令 $n\in\mathbb{N}_1$，我们要定义映射 $f:(\mathbb{Z},+,\cdot)\to(\mathbb{Z}_n,+,\cdot)$。对 $m\in\mathbb{Z}$，我们定义
\[
f(m)=m+n\mathbb{Z}
\]
则 $f$ 是一个环同态，而 $\ker(f)=n\mathbb{Z}\triangleleft R$。

加法的群同态已经证明过了，要证明的就是乘法幺半群同态。有乘法单位元，每一个环同态核都是理想，像都是子环。
\end{lemma}

\begin{theorem}[定理 2：环同态传递零元负元，交换环传递单位元]
假定 $R$ 与 $\bar{R}$ 是两个环，并且 $R$ 与 $\bar{R}$ 同态。那么，$R$ 的零元的像是 $\bar{R}$ 的零元；$R$ 的元 $a$ 的负元的像是 $a$ 的像的负元。并且，假如 $R$ 是交换环，那么 $\bar{R}$ 也是交换环；假如 $R$ 有单位元 $1$，那么 $\bar{R}$ 也有单位元 $\bar{1}$，而且 $\bar{1}$ 是 $1$ 的像。

我们要注意，一个环有没有零因子这一个性质经过了一个同态满射是不一定可以保持的。
\end{theorem}

\begin{proposition}[命题 2.8：$R$-$R'$ 环同态，核都是理想，像都是子环]
令 $f:(R,+,\cdot)\to(R',+,\cdot)$ 是一个环同态，则 $f$ 的核是 $R$ 的理想，$f$ 的像是 $R'$ 的子环。此即，
\begin{align}
\ker(f)&=\{a\in R:f(a)=0'\}\trianglelefteq R\\
\mathrm{im}(f)&=\{b\in R':\exists a\in R,\,b=f(a)\}=\{f(a)\in R':a\in R\}<R'
\end{align}
\end{proposition}

\begin{note}[定义2 9环同态]
加法和乘法都能拆括号（同态保持运算）；乘法单位元通过同态传递到像中。
\end{note}

\begin{theorem}[定理 3：同构传递整环、除环、域的性质]
假定 $R$ 同 $\bar{R}$ 是两个环，并且 $R\cong\bar{R}$。那么，若 $R$ 是整环，$\bar{R}$ 也是整环；$R$ 是除环，$\bar{R}$ 也是除环；$R$ 是域，$\bar{R}$ 也是域。

引理：假定在集合 $A$ 与 $\bar{A}$ 之间存在一个一一映射 $\varphi$，并且 $A$ 有加法和乘法。那么我们可以替 $\bar{A}$ 规定加法和乘法，使得 $A$ 与 $\bar{A}$ 对于一一对加法以及一一对乘法来说都同构。
\end{theorem}

\begin{theorem}[定理 4]
假定 $S$ 是环 $R$ 的一个子环，$S$ 在 $R$ 里的补集（就是 $R$ 中不属于 $S$ 的所有元作成的集合）与另一个环 $\bar{S}$ 没有共同元，并且 $S\cong\bar{S}$。那么存在一个与 $R$ 同构的环 $\bar{R}$，而且 $\bar{S}$ 是 $\bar{R}$ 的子环。
\end{theorem}

% ===================== 第三节 =====================
\section{理想与商环}

\begin{definition}[理想子环定义]
环 $R$ 的一个非空子集 $\mathfrak{U}$ 叫做一个\textbf{理想子环}，简称\textbf{理想}，是指
\begin{enumerate}[(i)]
  \item $a,b\in\mathfrak{U}\Rightarrow a-b\in\mathfrak{U}$；
  \item $a\in\mathfrak{U}$，$r\in R\Rightarrow ra,\,ar\in\mathfrak{U}$。
\end{enumerate}

以下两个理想是平凡的：
\begin{enumerate}
  \item 只包含零元的集合，这个理想叫做 $R$ 的\textbf{零理想}；
  \item $R$ 自己，这个理想叫做 $R$ 的\textbf{单位理想}。
\end{enumerate}
有的环除了这两个理想以外，没有其他理想，比方说除环。

任意 $a$ 属于理想，所以只要 $1$（单位元）在理想里面，这个理想就是整个环。理想是原环所有元素和 $a$ 相乘生成的元素组成的。

为什么要证明"某些东西加起来等于 $1$"？这通常是为了说明某个理想是整个环 $R$，即 $(a,b)=R$，我们要证明 $\exists r,s\in R,\;ra+sb=1$。因为只要 $ra+sb=1$，那么任意 $x\in R$ 都可以写成 $x=x\cdot 1=x(ra+sb)=(xr)a+(xs)b\in(a,b)$，所以说明 $(a,b)=R$。
\end{definition}

\begin{note}[定义2 10理想]
$(I,+)$ 是加法子群，左乘环中任意元素还在理想里面就是左理想，同理右理想。理想的符号 $\trianglelefteq$ 与正规子群一样，体现了理想与正规子群的类比。理想都是子环。
\end{note}

\begin{theorem}[定理 1：除环只有两个理想]
一个除环 $R$ 只有两个理想，就是零理想和单位理想。

\textbf{证明：} 假定 $\mathfrak{U}$ 是 $R$ 的一个理想而 $\mathfrak{U}$ 不是零理想。那么 $\mathfrak{U}\ni a\ne 0$，由理想的定义，$a^{-1}a=1\in\mathfrak{U}$，因而 $R$ 的任意元 $b=b\cdot 1\in\mathfrak{U}$。这就是说，$\mathfrak{U}=R$。证完。

因此，理想这个概念对于除环或域没有多大用处。
\end{theorem}

\begin{note}[定理1除环只有两个理想]
交换的除环就是域，所以域也只有两个理想：零理想和单位理想。这也是为什么理想这个概念对域没有什么作用。
\end{note}

\begin{example}[$\mathbb{Z}_n$ 剩余类环]
令 $R=\{\text{所有模 }n\text{ 的剩余类}\}$。我们替 $R$ 规定一种加法
\[
[a]+[b]=[a+b],
\]
并且知道 $R$ 对于这个加法来说成一个群。现在我们替 $R$ 规定一个乘法，我们规定
\[
[a][b]=[ab].
\]
模 $n$ 的剩余类是由整数间的等价关系 $a\equiv b\pmod{n}$ 当且仅当 $n\mid a-b$ 所决定的。如果 $[a]=[a']$，$[b]=[b']$，那么由等式 $ab-a'b'=a(b-b')+(a-a')b'$，容易证明 $[ab]=[a'b']$。由上述加法和乘法的定义易见：乘法适合结合律，并且两个分配律都成立。因此 $R$ 作成一个环。这个环叫做模 $n$ 的\textbf{剩余类环}。

如果 $n$ 不是素数，$n=ab$，$n\nmid a$，$n\nmid b$，那么在环 $R$ 里
\[
[a]\ne[0],\quad[b]\ne[0],\quad\text{但}\ [a][b]=[ab]=[n]=[0].
\]
因为 $[0]$ 正是 $R$ 的零元，这就是说，式 $[a][b]=0\Rightarrow[a]=0$ 或 $[b]=0$ 在 $R$ 里不成立。
\end{example}

\begin{proposition}[命题 2.6：交换环的理想，左理想就是右理想]
假设 $(R,+,\cdot)$ 是一个交换环，则 $I$ 是一个左理想当且仅当 $I$ 是一个右理想，又当且仅当 $I$ 是一个理想。

理想在乘法下吸收了整个环到理想上。
\end{proposition}

\begin{proposition}[命题 2.7：整环乘整数都是理想]
令 $n\in\mathbb{N}_1$，则 $n\mathbb{Z}\triangleleft\mathbb{Z}$，即
\[
n\mathbb{Z}\trianglelefteq\mathbb{Z}
\]
\textbf{证明：} 首先，我们知道 $(n\mathbb{Z},+)$ 是 $(\mathbb{Z},+)$ 的（加法）子群，这可以通过构建 $m\mapsto mn$ 的（加法）群同态来得到。其次，注意到 $\mathbb{Z}$ 是一个交换环，故我们只须验证 $R\mathfrak{U}\subset\mathfrak{U}$，在这里就是 $\mathbb{Z}\cdot n\mathbb{Z}\subset n\mathbb{Z}$。因为一个 $n$ 的倍数的数再乘任意数都还是 $n$ 的倍数。这促使我们，不要被看似简单的事实实现现象蒙蔽双眼，我们应当从中找到教学价值。
\end{proposition}

\begin{note}[理想并一下还是理想]
理想并一下还是理想。
\end{note}

\begin{definition}[定义 2.13：生成理想的定义]
令 $(R,+,\cdot)$ 是一个环，而 $A\subset R$。则 $(A)$，称为由 $A$ 生成的理想，定义为所有 $R$ 中包含 $A$ 的理想的交集，即
\[
(A)=\bigcap\{I\subset R:I\supset A,\;I\trianglelefteq R\}
\]
这里，我们用圆括号来区分尖括号，后者生成的是子环。所有包含他的理想的交集。生成的理想是包含 $(4)$ 的理想的交集。
\end{definition}

\begin{definition}[定义 2.14：主理想]
令 $(R,+,\cdot)$ 是一个环，而 $a\in R$，则我们定义
\[
(a)=(\{a\})
\]
称为由 $a$ 生成的\textbf{主理想}。一般地，若一个理想能被一个元素生成，我们就称其为主理想。对于 $a_1,\cdots,a_n\in R$，我们定义
\[
(a_1,\cdots,a_n)=(\{a_1,\cdots,a_n\})
\]
一般地，若一个理想能被有限个元素生成，我们就称其为有限生成的理想。

主理想要求就是由一个元生成的理想。证明是主理想就是证明一个元素能生成整个环；不是每个理想都能做到的，因为理想的诞生是理想中元素与环中元素作用并且不断补充进理想，而生成整个理想是理想里面的一个元素不与环作用。可交换+有单位元 $\Rightarrow$ 理想可以写成 $ra$ 形式，$r\in$ 环。
\end{definition}

\begin{definition}[剩余类环的定义与符号（即商环）]
$\bar{R}$ 叫做环 $R$ 的模 $\mathfrak{U}$ 的\textbf{剩余类环}。这个环我们用符号
\[
R/\mathfrak{U}
\]
来表示。
\end{definition}

\begin{note}[定义2 11商环]
$R$ 是一个环，$I$ 是 $R$ 的一个理想，商环就是 $R$ 中元素对 $I$ 做加法陪集形成的集合，$a+I$ 是一个集合。商群是加法的陪集，商环也是加法的陪集。类比商群的构造：正规子群对应理想，商群对应商环。
\end{note}

\begin{proposition}[命题 2.9：商环是环，当且仅当 $I$ 是理想]
令 $(R,+,\cdot)$ 是一个环，而 $I\trianglelefteq R$，则上述的加法和乘法是良定义的，且商环 $(R/I,+,\cdot)$ 是一个环。
\end{proposition}

\begin{theorem}[定理 1：理想和正规子群很像，模理想的剩余类还是环而且和原环同态]
假定 $R$ 是一个环，$\mathfrak{U}$ 是它的一个理想，$\bar{R}$ 是所有模 $\mathfrak{U}$ 的剩余类作成的集合。那么 $\bar{R}$ 本身也是一个环，并且 $R$ 与 $\bar{R}$ 同态。

\textbf{证明：} 映射 $a\longrightarrow[a]$ 显然是 $R$ 到 $\bar{R}$ 的一个同态满射，所以 $R$ 与 $\bar{R}$ 同态，而 $\bar{R}$ 是一个环。证完。
\end{theorem}

\begin{theorem}[定理 2：$R$ 和自己的商环同态，同态满射核就是理想]
假定 $R$ 同 $\bar{R}$ 是两个环，并且 $R$ 与 $\bar{R}$ 同态，那么这个同态满射的核 $\mathfrak{U}$ 是 $R$ 的一个理想，并且
\[
R/\mathfrak{U}\cong\bar{R}.
\]
核在环里面是映射到 $0$ 元。
\end{theorem}

\begin{theorem}[定理 3：$R$ 和 $R$ 的商环同态满射的性质，传递]
在环 $R$ 到环 $\bar{R}$ 的一个同态满射之下：
\begin{enumerate}[(i)]
  \item $R$ 的一个子环 $S$ 的像 $\bar{S}$ 是 $\bar{R}$ 的一个子环；
  \item $R$ 的一个理想 $\mathfrak{U}$ 的像 $\bar{\mathfrak{U}}$ 是 $\bar{R}$ 的一个理想；
  \item $\bar{R}$ 的一个子环 $\bar{S}$ 的原像 $S$ 是 $R$ 的一个子环；
  \item $\bar{R}$ 的一个理想 $\bar{\mathfrak{U}}$ 的原像 $\mathfrak{U}$ 是 $R$ 的一个理想。
\end{enumerate}
\end{theorem}

% ===================== 第四节 =====================
\section{主理想、整环与极大理想}

\begin{definition}[左零因子、右零因子]
如果在一个环里
\[
a\ne 0,\quad b\ne 0,\quad\text{但}\;ab=0,
\]
我们就说，$a$ 是这个环的一个\textbf{左零因子}，$b$ 是一个\textbf{右零因子}。左、右零因子都称作\textbf{零因子}。交换环左零因子就是右，非交换环不一定。
\end{definition}

\begin{proposition}[命题 2.10：由子集生成的理想仍是理想]
令 $(R,+,\cdot)$ 是一个环，而 $A\subset R$，则 $(A)\trianglelefteq R$。
\end{proposition}

\begin{proposition}[命题 2.11：交换环中主理想的表达式]
令 $(R,+,\cdot)$ 是一个交换环，而 $a\in R$，则
\[
(a)=Ra=\{ra:r\in R\}.
\]
一般地，若 $a_1,\cdots,a_n\in R$，则
\[
(a_1,\cdots,a_n)=Ra_1+\cdots+Ra_n=\{r_1a_1+\cdots r_na_n:r_1,\cdots,r_n\in R\}.
\]
\end{proposition}

\begin{definition}[定义 2.18：整环没有零因子是交换环]
令 $(R,+,\cdot)$ 是一个环，则我们称 $R$ 是个\textbf{整环}，若它是个非零交换环，且没有零因子，即
\begin{align}
R&\ne\{0\}\\
&R\text{ 是个交换环}\\
\forall a,b\in R,\;(ab=0&\Longrightarrow a=0\;\text{或}\;b=0)
\end{align}
若 $a\ne 0$ 满足 $\exists b\ne 0$ 使得 $ab=0$，我们就称其为一个零因子。特别地，类加群的时候，凡事不和阶互素的元都是零因子。
\end{definition}

\begin{note}[零因子与消去律互相证明]
无零因子 $\Leftrightarrow$ 消去律成立。
\end{note}

\begin{definition}[特征的定义，对加法所有非零元素共同的阶]
一个无零因子环 $R$ 的非零元的相同的（对加法来说的）阶叫做环 $R$ 的\textbf{特征}。

这样，一个没有零因子的环 $R$ 的特征如果是无限的，那么在 $R$ 里计算规则 (1) 式永远是对的；$R$ 的特征如果是有限数，那么这个计算规则就永远不对。特征是一个很重要的概念，因为它对环和域的构造都有决定性的作用。
\end{definition}

\begin{note}[定理2 整环除环域的特征]
特征是加法中所有非零元素共同的阶。整环、除环、域都没有零因子，可以证明特征只能是素数或无限（不能是合数）。域一定有特征，有限域的特征只能是素数 $p$。
\end{note}

\begin{theorem}[定理 1：无零因子环所有不为 $0$ 的元对加法阶都是一样的]
在一个没有零因子的环 $R$ 里所有不等于零的元对于加法来说的阶都是一样的。
\end{theorem}

\begin{definition}[主理想环的定义]
一个整环 $I$ 叫做一个\textbf{主理想环}，是指 $I$ 的每一个理想都是一个主理想。

我们说，一个主理想环一定是一个唯一分解环。为证明这一点，我们需要两个引理，这两个引理本身也很重要。
\end{definition}

\begin{note}[定义2 25主理想整环]
主理想整环：每一个理想都是主理想（由单个元素生成）。主理想整环一定是唯一分解环；不可约元可以生成极大理想。
\end{note}

\begin{definition}[极大理想的定义]
一个环 $R$ 的一个不等于 $R$ 的理想 $\mathfrak{U}$ 叫做一个\textbf{极大理想}，是指除了 $R$ 同 $\mathfrak{U}$ 自己以外，没有包含 $\mathfrak{U}$ 的理想。

不等于父环；第二，不被其他理想包含。两步走：$I\ne R$，包含 $I$ 的理想只有 $R$。
\end{definition}

\begin{lemma}[引理 1：极大理想对剩余类环充要]
假定 $\mathfrak{U}\ne R$ 是环 $R$ 的理想。剩余类环 $R/\mathfrak{U}$ 除了零理想同单位理想以外不再有其他的理想，当且仅当 $\mathfrak{U}$ 是极大理想时。

商环只有零理想和单位理想，当被除的理想是极大理想。
\end{lemma}

\begin{note}[极大理想引理1对主理想整环的推论]
在主理想整环中，一个主理想 $(a)$ 是极大理想当且仅当 $a$ 是不可约元素。
\end{note}

\begin{lemma}[引理 2：交换环只有零理想和单位理想 $R$ 是一个域]
如果有单位元（$\ne 0$）的交换环 $R$ 除了零理想同单位理想以外没有其他的理想，那么 $R$ 一定是一个域。

交换环+只有零理想和单位理想就是一个域。
\end{lemma}

\begin{theorem}[极大理想与域]
假定 $R$ 是一个有单位元的交换环，$I$ 是 $R$ 的一个理想，则
\[
R/I\text{ 是一个域}\iff I\text{ 是极大理想}.
\]
\end{theorem}

% ===================== 第五节 =====================
\section{素理想与整数环中的例子}

\begin{lemma}[引理 2.7：素数的性质]
若 $p$ 是一个素数，$a,b\in\mathbb{Z}$，则
\[
p\mid ab\iff p\mid a\;\text{或}\;p\mid b
\]
注意到这个命题对于合数 $n$ 是错的。
\end{lemma}

\begin{lemma}[引理 2.8：合数的性质]
若 $n$ 是一个合数，则存在 $a,b\in\mathbb{Z}$，使得
\[
n\mid ab,\quad n\nmid a,\quad n\nmid b
\]
\textbf{证明：} 证明是简单的，若 $n$ 是一个合数，我们可以写一个平平凡分解 $n=ab$，其中 $a,b\ne\pm1$，则 $n\mid ab$，可是 $|a|<n$ 故 $n\nmid a$，同理 $n\nmid b$。

由此，在整数中，这个性质只对素数成立的。实际上，如果我们用理想来看，这个性质也可以改为
\[
\forall a,b\in\mathbb{Z},\;(ab\in p\mathbb{Z}\iff a\in p\mathbb{Z}\;\text{或}\;b\in p\mathbb{Z})
\]
因此，在一般的交换环中，我们定义素理想就是满足这个性质（如果两个元素的积在理想中，则至少有一元素在理想中）的理想，但不是每一个元素都是零的环的理想。
\end{lemma}

\begin{example}[整数环中，素数生成的都是极大理想]
\textbf{例 1} \quad 我们看整数环 $R$。我们说，由一个素数 $p$ 所生成的主理想 $(p)$ 是一个极大理想。因为假定 $\mathfrak{B}$ 是一个不等于 $(p)$ 的 $R$ 的理想，并且
\[
\mathfrak{B}\supseteq(p),
\]
那么 $\mathfrak{B}$ 一定包含一个不能被 $p$ 整除的整数 $q$。由于 $p$ 是素数，$q$ 与 $p$ 互素，所以存在整数 $s,t$，有
\[
sp+tq=1.
\]
包含一个和 $p$ 互素是这个理想包含由 $p$ 生成的理想的必要条件，也就是存在 $p$ 生成不了的元，然后用互素条件就能把关键的 $1$ 证明在这个包含 $p$ 的理想里面，有 $1$ 的理想就是整个环。
\end{example}
